% Preamble-tillegg for LaTeX-eksport av prosjektrapporten (inkluderes via pandoc -H)
% Luftigere kapitteloverskrifter
\usepackage{titlesec}
\titleformat{\section}{\Large\bfseries}{\thesection}{0.6em}{}
\titleformat{\subsection}{\large\bfseries}{\thesubsection}{0.5em}{}

% Bilde- og tabelltekst i kursiv og litt mindre, midtstilt (jf. typografireglene)
\usepackage{caption}
\captionsetup{font={small,it},labelfont={small,it},justification=centering}

% Midtstill longtable-tabeller horisontalt på siden.
\ifdefined\LTleft
  \setlength{\LTleft}{0pt plus 1fill}
  \setlength{\LTright}{0pt plus 1fill}
\fi

% Avsnitt med luft og uten innrykk
\setlength{\parskip}{0.4em}
\setlength{\parindent}{0pt}

% Slå av LaTeX sin automatiske seksjonsnummerering, siden kapitlene
% allerede har manuelle numre i overskriftene (f.eks. «1. Innledning»).
% Hindrer dobbel nummerering som «1 1. Innledning».
\setcounter{secnumdepth}{0}

% Innholdsfortegnelsen viser kapitler og delkapitler (ikke underdelkapitler),
% slik at den blir oversiktlig.
\setcounter{tocdepth}{2}

% Lær pdflatex de spesialtegnene rapporten bruker i løpende tekst, slik at
% kompilering ikke stopper på «Unicode character not set up for use with LaTeX».
% Beskyttet med \ifdefined slik at det er trygt også under xelatex/lualatex
% (der tegnene håndteres direkte og \DeclareUnicodeCharacter ikke finnes).
\ifdefined\DeclareUnicodeCharacter
  \DeclareUnicodeCharacter{2212}{\textminus}            % − ekte minustegn
  \DeclareUnicodeCharacter{2192}{\textrightarrow}       % → høyrepil
  \DeclareUnicodeCharacter{2264}{\ensuremath{\leq}}     % ≤ mindre enn eller lik
  \DeclareUnicodeCharacter{2248}{\ensuremath{\approx}}  % ≈ omtrent lik
  \DeclareUnicodeCharacter{00D7}{\texttimes}            % × gangetegn
  \DeclareUnicodeCharacter{00B2}{\textsuperscript{2}}   % ² opphøyd to (CV²)
\fi
